\begin{center}
	\LARGE
	Bake a Cake
\end{center}

For this assignment you're baking a cake. It is in groups of your own choosing. The amount of each ingredient is found by solving a series of problems. The recipy is for a carrot cake with cream cheese frosting stolen from a not specifically disclosed but quite widespead Danish food blogger.

\textbf{Ingredients}
\begin{itemize}
	\item[$\cdot$] \underline{\phantom{150}} g butter at room temperature.
	\item[$\cdot$] \underline{\phantom{100}} g sugar.
	\item[$\cdot$] \underline{\phantom{100}} g brown sugar (brun farin).
	\item[$\cdot$] \underline{\phantom{4}} eggs.
	\item[$\cdot$] \underline{\phantom{200}} g wheat flour.
	\item[$\cdot$] \underline{\phantom{2}} tea spoons baking powder.
	\item[$\cdot$] \underline{\phantom{1}} tea spoons ground cinnamon.
	\item[$\cdot$] \underline{\phantom{0.5}} tea spoons salt.
	\item[$\cdot$] \underline{\phantom{275}} g peeled and coarsely grated carrots.
\end{itemize}

\textbf{Ingredienser til topping}
\begin{itemize}
	\item[$\cdot$] \underline{\phantom{125}} g butter at room temperature.
	\item[$\cdot$] \underline{\phantom{1}} tea spoon vanilla sugar 
	\item[$\cdot$] \underline{\phantom{125}} g powdered sugar.
	\item[$\cdot$] \underline{\phantom{125}} g natural cream cheese.
	\item[$\cdot$] Peel and juice from \underline{\phantom{1}} organic lemon.
\end{itemize}

\textbf{Procedure}
Whip butter, sugar and brown sugar till airy. Whip in one egg at a time such that the mixture won't split.

Mix the dry ingredients in another bowl and fold in the dry ingredients into the batter. At last, fold in the grated carrots and bake the cake in a 20cm springform pan at 170$^\circ$C convection oven for about 45 minutes. Make sure to grease the pan well before adding in the batter.

\textbf{Procedure for frosting}
Whip butter, sugar and vanilla sugar until white and airy. Mix in the lemon juice, lemon peel and cream cheese and add it on the cake when the cake is completely cold. 

\newpage

\begin{opgavetekst}{Problem 1}
\end{opgavetekst}	
\begin{delopgave}{1}
	\begin{align*}
		\frac{601}{4} + \frac{52}{2} - \frac{210}{8}
	\end{align*}
\end{delopgave}
%
%
\begin{opgavetekst}{Problem 2}
	A polynomium $p$ is given by
	\begin{align*}
		p(x) = ax^2 - 500x + 600.
	\end{align*}
\end{opgavetekst}
\begin{delopgave}{1}
	Determine $a$ such that $p$ has roots in $x = 2$ and $x = 3$.  
\end{delopgave}
\begin{meretekst}
	Another polynomium $q$ is given by
	\begin{align*}
		q(x) = x^2 + 99x + c.
	\end{align*}
\end{meretekst}
\begin{delopgave}{2}
	Determine $c$ such that the graphs of $p$ and $q$ intersect when $x = 1$. 
\end{delopgave}
%
%
\begin{opgavetekst}{Problem 3}
	Two lines $l$ and $m$ are given by
	\begin{align*}
		&l: \ -4x + 2y = -20 \\
		&m: \ 5x + y = 18
	\end{align*}
\end{opgavetekst}
\begin{delopgave}{1}
	Determine the $x$-coordinate for the intersection between $l$ and $m$. 
\end{delopgave}
\begin{delopgave}{2}
	Determine the $y$-coordinate for $l$ when $x = 105$. 
\end{delopgave}
%
%
\begin{opgavetekst}{Problem 4}
	A function $f$ is given by
	\begin{align*}
		f(x) = \ln(2x-3).
	\end{align*}
\end{opgavetekst}
\begin{delopgave}{1}
	Determine $f^{-1}(0)$
\end{delopgave}
%
%
\begin{opgavetekst}{Problem 5}
	An exponential function $g$ is given by
	\begin{align*}
		g(x) = 7\cdot 1.01^x
	\end{align*}
\end{opgavetekst}
\begin{delopgave}{1}
	Determine how many $\%$ $g$ grows by if $x$ is increased by 1. 
\end{delopgave}
%
%
\begin{opgavetekst}{Problem 6}
	A polynomium $q$ is given by
	\begin{align*}
		q(x) = ax^2 + 5x + 12.5
	\end{align*}
\end{opgavetekst}
\begin{delopgave}{1}
	Determine $a$ such that $q$ has exactly one root. 
\end{delopgave}
%
%
%
%
\begin{opgavetekst}{Problem 7}
	A sequence $\{a_i\}$ is given by
	\begin{align*}
		a_i =
		\begin{cases}
			25  &\textnormal{ if }i = 1,\\
			a_{i-1} + 50 &\textnormal{ if } i > 1.
		\end{cases}
	\end{align*}
\end{opgavetekst}
\begin{delopgave}{1}
	Determine $a_6$.
\end{delopgave}
\begin{delopgave}{2}
	Determine $i$ such that $a_i = 6225$ 
\end{delopgave}
%
%
\begin{opgavetekst}{Problem 8}
	A function $f$ is given by
	\begin{align*}
		f(x) = \sqrt{5x}.
	\end{align*}
\end{opgavetekst}
\begin{delopgave}{1}
	How far should the graph of $f$ be shifted to the left to pass through $(4,5)$? 
\end{delopgave}
%
%
\begin{opgavetekst}{Problem 9}
	You're given a yearly interest of 2.8$\%$ in a bank account. You deposit $100 \TS 000$ kr in the bank. 
\end{opgavetekst}
\begin{delopgave}{1}
	How much money is in the account after 8 years. The result should be given in 1000 kr. 
\end{delopgave}
\begin{delopgave}{2}
	Determine the number of years before $3 \TS 156 \TS 016$ kr is the balance of the account.
\end{delopgave}
%
%
\begin{opgavetekst}{Problem 10}
	A piecewise defined function $g$ is given by
	\begin{align*}
		g(x) = 
		\begin{cases}
			x^2 + 4x - 3,  &\textnormal{ if }x \leq 2,\\
			4x + b,  &\textnormal{ if }x > 2.
		\end{cases}
	\end{align*}
\end{opgavetekst}
\begin{delopgave}{1}
	Determine the number $b$ such that $g$ is continuous in $x = 2$. 
\end{delopgave}
